\section{Conclusion}

Chunk length is usually treated as a per-task constant, and the quantity it
trades against, the rate at which error compounds while the policy acts blind,
is usually not measured at all. We measured it directly as a finite-time
maximal Lyapunov exponent per state, with a procedure that trains nothing and
reads no policy internals.

On Claim 1 the evidence is strong. Across more than twenty datasets and three
platforms, both stability regimes occur inside almost every episode, with
several crossings each. The unstable segments are a small minority of
timesteps, short, and concentrated at contact. The structure survives a
fourfold change in probe size. Each task's response to fixed chunk length
follows its label profile, and this holds for two different policy classes.

On Claim 2 the evidence is also strong. Small heads reading only camera frames
and proprioception recover the regime at 0.83 to 0.96 AUROC on confident stamps
in domain, transfer across tasks at 0.72 to 0.77, and beat a privileged contact
baseline by a wide margin. Temporal history does not help, which places the
signal in the current configuration rather than in motion. Audited against our
labels, the replan signals used by existing adaptive-horizon methods do not
track the measured stability, and one of them is anti-correlated with it. The
closed-loop axis is measured for four robomimic tasks and in progress on
RoboCasa.

On Claim 3 the evidence is negative so far, and we prefer to report that
plainly. Learned predictors driving the switch gain a few points over a fixed
16-step chunk, never collapse the way constant single-step execution does, but
stay inside the run-to-run spread and never beat a properly swept constant
chunk. A paired fork test, which removes episode variance by running both
decisions from one simulator state, finds that closing the loop at a measured
unstable moment rescues as many episodes as it breaks. A 42-cell sweep with
diffusion policies on three other benchmarks agrees. Along the way we found two
deployment failures that offline detection metrics hide: a regression head
compresses its predictions so the label threshold never fires, and a threshold
calibrated on demonstrations mis-fires by up to a factor of seven once the
policy visits its own states.

The open question is the second stability axis. Switching should pay exactly
where a segment is open-loop unstable and closed-loop stable, and we cannot yet
say how common that combination is, because a closed-loop measurement needs
both branches under policy control and a control for the policy's own sampling
noise. That measurement, milder switches than single-step execution, and paired
tests on more tasks are the direct next steps.
